\documentclass[a4paper]{article}

\usepackage{graphicx}
\usepackage{amsmath,amssymb}
\usepackage{minted}
\usepackage{multirow}
\usepackage{float}
\usepackage{todonotes}
\usepackage{forest}
\usepackage{xstring}
\usepackage{xcolor}
\usepackage[most]{tcolorbox}
\usepackage{xparse}
\usepackage{natbib}

\usetikzlibrary{graphs,graphdrawing}
\usegdlibrary{layered}

% for external graphviz
\input{/home/donkarlo/Dropbox/repo/nd_language_ai_project/src/nd_language_ai/formal/computer/programming/tex/latex/code_clips/external_graphviz.tex}

% for tags
\input{/home/donkarlo/Dropbox/repo/nd_language_ai_project/src/nd_language_ai/formal/computer/programming/tex/latex/parts/control_sequence/kind/semtag/semtag.tex}

% for links
\input{/home/donkarlo/Dropbox/repo/nd_language_ai_project/src/nd_language_ai/formal/computer/programming/tex/latex/code_clips/seehere.tex}

\title{Refractory period}
\date{}
\author{Mohammad Rahmani}

\begin{document}
    \maketitle
    Refractoriness is the fundamental property of any object of autowave nature (especially excitable medium) not responding to stimuli, if the object stays in the specific refractory state. In common sense, refractory period is the characteristic recovery time \seeherefooter{https://en.wikipedia.org/wiki/Refractory_period_(physiology)}.
    
    
    \section{Usage ofthis period}
        \begin{enumerate}
            \item It prevents the neuron from firing action potentials continuously without a time interval.
            \item It limits the maximum firing rate of the neuron.
            \item It keeps consecutive action potentials temporally separated.
            \item It allows the ion channels to return to their ready state after each action potential.
        \end{enumerate}
        \seeherefooter{https://en.wikipedia.org/wiki/Refractory_period_(physiology)#Neuronal_refractory_period}.
%    \bibliographystyle{plainnat}
%    \bibliography{/home/donkarlo/Dropbox/repo/nd_shared_research/bibliography/bibliography.bib}
\end{document}